\scanpage{6}
En termes des bases (8) (8$'$), $\varphi$ sera donnée par une matrice (notée \deleted{\unsure{aussi} $\varphi$ \unsure{$\underline{\mathrm O}_S$}})
\[
 \varphi=(\varphi_{ij})_{1\leq i,j\leq3}\qquad \varphi_{ij}=\varphi(e_i,e_j) \tag*{(11)}
\]
-- la matrice de $f$ est alors \uline{transposée} de la matrice $\varphi$ -- et on trouve que la structure d'algèbre alternée correspondante est donnée par
\[
 [e'_2\unsure{\wedge|,}e'_3]_\varphi=\varphi\alpha'(e'_2\wedge e'_3)=\varphi\alpha'(\bar e'_1)
 \underset{(10)}{=}\varphi(e_1)=\sum_j\varphi_{1j}e'_j,
\]
et de \deleted{$[e'_3\wedge e'_1]=$} même pour les deux autres \deleted{produits}\inserted{crochets}
\[
\left\{\begin{aligned}
 [e'_2,e'_3]_\varphi&=\sum_j\varphi_{1j}e'_j\\
 [e'_3,e'_1]_\varphi&=\sum_j\varphi_{2j}e'_j\\
 [e'_1,e'_2]_\varphi&=\sum_j\varphi_{3j}e'_j
\end{aligned}\right. \tag*{(12)}
\]
Ceci posé, on trouve (laissant tomber l'indice $\varphi$ dans les \inserted{crochets})
\[
 [[e'_1\wedge e'_2],e'_3]=\sum_j(\varphi_{32}\varphi_{1j}-\varphi_{31}\varphi_{2j})e'_j \tag*{(13)}
\]
et de même pour les expressions \unsure{déduite|déduites} par permutations circulaires
\[
\left\{\begin{aligned}
 [[e'_2\wedge e'_3],e'_1]&=\sum_j(\varphi_{13}\varphi_{2j}-\varphi_{12}\varphi_{3j})e'_j\\
 [[e'_3\wedge e'_1],e'_2]&=\sum_j(\varphi_{21}\varphi_{3j}-\varphi_{23}\varphi_{1j})e'_j
\end{aligned}\right. \tag*{(13 bis)}
\]
Les \deleted{coefficients} symétries des formules apparaissent mieux en posant
\[
 \deleted{\varphi_{11}=a,\ \unsure{\varphi_{1\cdot}|\varphi}=}
\]
\[
\left\{\begin{aligned}
 \varphi_{11}=a,\quad &\varphi_{22}=b,\quad \varphi_{33}=c\\
 \varphi_{23}=p,\quad &\varphi_{31}=q,\quad \varphi_{12}=r\\
 \varphi_{32}=p',\quad &\varphi_{13}=q',\quad \varphi_{21}=r'
\end{aligned}\right. \tag*{(14)}
\]
i.e. en écrivant
\[
 \varphi=\begin{pmatrix}a&r&q'\\ r'&b&p\\ q&p'&c\end{pmatrix} \tag*{(15)}
\]
de sorte que la matrice transposée de $\varphi$ s'obtient en échangeant $p$ avec $p'$, $q$ avec $q'$, $r$ avec $r'$, et

% archive sheet number: 5
